\chapter{Alpha As A Bracket}

\section{The coupling is asked}

The field closed and posed its question. This is the chapter that answers --- and it opens not with a number but with the
\emph{shape} of the answer, because the shape is the first surprising thing about it. When the instrument asks for the
coupling, the coupling does not return as a magnitude lying out in the world, waiting to be read off a scale. It returns as a
\emph{self-energy} --- a cost the apparatus pays, in its own pulse, to weigh its own account of the electron. The coupling is
asked, and answered, as the self-energy the last part built.

Take the plain construction first, because beneath the physical name there is an exact and modest claim. The apparatus holds
two descriptions of one object: the electron as a \emph{source}, a charge that sits and sources a field, and the electron as
a \emph{field}, the configuration that charge establishes around itself. These are two accounts of a single thing, and if the
apparatus's bookkeeping were perfect they would cancel --- the source paid for exactly by the field it makes. They do not
quite cancel. Weighed against each other in the apparatus's own currency, up to the fixed calibration the construction
carries between the two scales, a remainder survives: the two descriptions fail to cancel, and the surviving remainder is the
residue. That residue is the coupling. To ask the instrument for the coupling is to ask what it costs the instrument to weigh
its two accounts of the electron and find they do not close.

The self-energy is exactly this remainder read as a physical quantity, and the reading is marked. That the coupling is
\emph{asked as a residual} --- the leftover when the apparatus weighs its source-account against its field-account, up to the
construction's own calibration --- is the construction's, built at the last chapter and pointed here. That this residual
\emph{is} the self-energy, the quantity physics knows by that name for the field a charge carries with itself, is the book's
reading laid over the residual, and it is marked as a reading: the instrument computes a remainder; the identification of that
remainder with the physical self-energy is the interpretation, not an equation proved in these pages. The concept that the
apparatus is naming its own account --- extrapolating its two descriptions toward the value they fail to share --- is carried
here as an idea only, not as a derivation the book leans on.

And the currency is the apparatus's own, exactly as the last part insisted. The residual is a count of the apparatus's own
pulse, scaled by the construction's own fixed calibration into its own units --- no external standard anywhere in it. So the
coupling, read as a self-energy, is denominated as every reading in this book has been: the instrument's own pulse, over its
own constants, spent on its own account of the electron. It reaches outside itself for nothing, not even here, at the number
it was built to find.

And here, opening the payoff, the chapter does what every chapter before it has done at this edge. It names the kind and
withholds the value. It has said \emph{what} the coupling is: a residual, a self-energy, the cost of the apparatus weighing
its two accounts of the electron. It has not said, and will not say here, \emph{how much}. The magnitude is two sections off;
and even there it will not arrive as a value the instrument derived, but as an external number the instrument \emph{proves it
cannot compute}, set beside the bracket the instrument can. In the two verbs, the asking funges the apparatus's two accounts
into one residual --- bags them, weighs them, names the leftover ``the coupling'' --- and refuses, this section, the last
tange that would price the bag. The kind is in hand; the number is one section ahead.

\section{The cut, come to rest at the count}

The last section named what the coupling is --- a residual, a self-energy --- and kept its number to itself. This section
takes the number. It is the first place in the book where the instrument reads a magnitude off its own working and writes it
down, and so it is where the book's whole discipline about numbers is finally tested against an actual reading. The reading
is not a point. It is a bracket, and the width of the bracket is the honest part.

Recall how the instrument reaches a value it cannot compute exactly: it brackets and refines. But it refines on its own
arithmetic --- the plain ratios it is built from --- and not on the real line laid over them. It walls the quantity between a
lower ratio and an upper, and closes the gap by the \emph{mediant}: from a lower bound $p/q$ and an upper $r/s$ it forms the
in-between ratio $(p+r)/(q+s)$, a step of pure integer arithmetic that never leaves the rationals and borrows no continuum to
take it. Iterated, the mediant generates the \emph{convergents} of the value the instrument is closing on --- its successive
best rational approximations --- and it would generate them forever if nothing stopped it. Something stops it: the resolution
floor of Chapter 6, the depth below which the thermal jitter of the reading frame erases the distinction between two
candidates. Here that floor has a name taken from the number of refinements the instrument can meaningfully make before it
hits bottom --- it counts to three, and no further. The refinement comes to rest at the count of three. It does not refine into
noise; it stops where its resolution stops, and it can prove the bracket it stops on is genuine --- that its lower wall lies
below its upper --- by direct decision on its own computed walls.

So the instrument, asked for the inverse of the coupling, hands back not a number but two: a lower wall and an upper wall,
the value somewhere between them and no obligation to say where. Read off the instrument's own run, at its own count-to-three
resolution, the two walls are the two convergents that three refinements leave standing --- and, read as the inverse of the
coupling, the lower wall sits near $129.6$ and the upper near $137.7$. These are the instrument's own readings, in its own
units, computed by its own construction --- labeled, in the construction's own words, ``not a magnitude derivation.'' They
are where the refinement comes to rest, not a value the instrument claims to have pinned. Two honesties belong with these two
numbers. The coming to rest --- the stopping at three --- is the robust thing, a property of the resolution and not a tuning. And the
\emph{width}, which is wide, is honest: the lower wall is loose because three refinements leave the coarse convergent
standing there, and the instrument claims no finer bound than three counts have earned. Nothing about that width is brought
in from outside: the mediant generates its own subdivision from the instrument's own ratios, with nothing chosen
beforehand. The bracket is the instrument's honest resolution at three counts, self-generated and nothing more.

And here is why the count is \emph{three} rather than any other depth, which is the deepest content of the section. The value
the instrument closes on is not arbitrary: the second variation gives the coupling an inverse-square law, the residual
falling off as the square of a distance, so the crossing --- where the residual meets the instrument's own target --- sits at
an exact square root, $\sqrt{18/5}$, a ratio of the instrument's own constants under a root sign. That is a quadratic surd,
and a quadratic surd has a continued fraction that \emph{repeats}: its partial quotients fall into a fixed cycle, by a
theorem of Lagrange's. So the instrument's own subdivision of the crossing is periodic and self-generating, closing on a
value it draws from its own numbers with nothing brought in from outside. The count-to-three is exactly three of that
continued fraction's partial quotients --- three levels of the self-generating subdivision --- and three levels place the
crossing between the two convergents that wall the wide bracket above. That is the ``why three'' the bracket rides on: three
levels of the instrument's own periodic subdivision, and no more, because the resolution floor allows no more.

There is a \emph{second} ``three'' in the reading, and the two must not be confused. To read the residual the instrument also
integrates the second variation, and the natural rule there is a three-node one --- a Gaussian quadrature with nodes near the
roots of the degree-three Legendre polynomial, exact in principle to the fifth degree. But that three is a matter of how the
integrand is \emph{sampled}, a side reading, and it rests on rational approximations of those roots --- it is carried as a
marked convenience, not a derived fact, and it is \emph{not} the reason the bracket has the width it has. The bracket's three
is the resolution-depth three: three partial quotients of the crossing's own continued fraction. The integration's three is
separate and softer, and the book keeps them apart.

And here is why the instrument \emph{stops} at three rather than pressing further. You could take more convergents past the
floor. The walls would keep closing, and the continued fraction would eventually settle onto a single value. But that value
would be a fiction of over-refinement. Past the resolution floor the subdivision is no longer resolving anything real --- it
is ringing, converging smoothly and convincingly onto a value of the instrument's \emph{own} --- near $137.011$ --- the
number its own arithmetic settles toward when it refines past the floor, a value that looks like an answer and is an artifact
of refining past what the counting earns. An instrument that keeps refining past its floor produces a crank's point: precise,
confident, and wrong. The instrument that comes to rest at the count produces a bracket: wider, humbler, and honest. The strain the whole book has tracked --- the pressure to
force a value where the floor allows only a bound --- is here relaxed, deliberately, into the width of the open interval. In
the two verbs, the refinement funges the coupling into the bracket and then, at the floor, declines to tange it finer: below
the count-to-three floor no characteristic selects one candidate from another, so the instrument funges the value into the
interval and leaves the last selection un-made. What this section does \emph{not} do --- and the next makes the reason
exact --- is claim the bracket contains the world's value. That is a further question, and the honest answer is not the one a
crank would give.

There is a sharper reason the width is honest, and it lives in the law itself. The inverse-square field runs to infinity
at one point --- a pole, where the map from a crossing-distance to the coupling divides by zero and holds no finite
value --- and the crossing $\sqrt{18/5}$ the descent reaches for sits just past it, so every bisection that closes toward
the value straddles that singular point. A hard cut taken across such a point does not come back clean; it rings, and the
width is that ring settled --- the \emph{stabilizer} of an unstable point that holds no value of its own, not a coarseness
the instrument failed to sharpen. In the physical reading the ring is a ringdown, the tone a struck and unstable thing
gives off as it settles; but the reading is a lens and no more, marked hard --- the electron is not two masses merging,
and the ringdown only \emph{reads like} the width, a resemblance and never a shared process, the no-name crossing the
construction has refused throughout.

\section{The bracket, not the real}

The last section gave the instrument's honest reading of the coupling: a bracket, computed off its own run, come to rest at its own
floor, held open. This closing section does the one thing the whole book has been building toward and holding back --- it
names the external value, the number a laboratory measures for this coupling out in the world. It names it exactly once, and
names it as what it is: external. And having named it, it proves the thing the whole discipline has been for --- that the
instrument cannot derive it.

There is a number for this coupling in the world, measured by physicists in an apparatus of their own: a single electron held
in a trap, its behavior read off to great precision. That measured value --- the inverse of the coupling --- the laboratories
put at $137.036$~\cite{hanneke2008}. This is the one place in the entire volume that number appears. It has not been used
before --- not smuggled into a definition, not fed to the instrument as a target, not tuned toward. The instrument never
fetched it; nothing in the construction reached out to a table of constants and read it in. It is named here, once, for the
first and only time, and strictly as a fact about the world's laboratories, not a fact the instrument produced. The
instrument's own reading is the bracket of the last section; this number is the world's; and the two are kept rigorously
apart, never set against each other in any gesture of nearness.

And here is the claim the whole book was built to earn, and it is not the claim a crank would make. A crank, holding a
bracket and a lab value, would announce that the instrument had predicted the value. This book does the opposite. Its claim is
that the instrument \emph{cannot derive} the external value --- and that this is not a shortcoming but a theorem. The
magnitude the laboratory measures lives in the part of the object that counting cannot reach: it is invariant under exactly
the operations a finite instrument performs, so no amount of the instrument's counting, however deep, arrives at it. The
instrument can prove this about itself --- that what it does, count and bisect and come to rest, is structurally unable to produce
that particular number. So what the instrument offers is not a derivation of the lab value and not a prediction of it. It is
a bracket of its own, floored at its own resolution, and a proof that the world's value is not something the instrument's
kind of work can compute. The number is named, the inability to reach it is proved, and those are the same act of honesty.

The book will not compare the two, and the refusal is deliberate and total. It does not report where the external number
falls relative to the instrument's walls, does not measure the distance between them, does not cash any nearness for credit.
There are two exact reasons. The first: the instrument's own reading is not a point but a wide bracket, and even the value
its own arithmetic settles toward, pressed past the floor, is a number of \emph{its own} --- not the world's, and never meant
to be. The instrument settles on a value of its own; the laboratory measures another; that these are different numbers is not
a defect to be closed but part of the very fact the section proves. The second is deeper, and is the
whole point: the instrument's provable claim is the \emph{reverse} of a capture. It does not say ``the value is inside my
walls, so I found it.'' It says ``the value is a thing I can prove my counting cannot reach.'' To set the two numbers side by
side and remark on their nearness would be to unsay exactly the theorem the book exists to prove. The reading the instrument
stands behind is the bracket. The real value is the world's, named once, and left where it lives.

In the two verbs, the instrument funges the coupling into its own bracket --- bags it, owns it, at its own floor --- and
refuses, at the last, to tange the world's point out of that bag. It cannot select the external value; it has proved it
cannot count its way to it; so it does not pretend to have. The book closes its payoff chapter exactly where its honesty
requires: with a bracket it owns and a number it does not, the two named in their separate places and never confused. The
coupling, read by this instrument, is the bracket, floored and open. The coupling, measured in the world, is a number the
instrument can name and prove it cannot derive. That those are two different things, honestly kept apart, is the whole of
what the physical gauge was built to say --- and the jar the last part opens is nothing more than this bracket, held out
plainly, with the world's number set honestly apart from it and never claimed as its own.
